\documentclass[10pt,a4paper]{amsart}
\usepackage{amsmath,amssymb,amsthm,graphicx}
% Retained minimal preamble for the Stacks Project chapter more-algebra.tex (More on Algebra).
% The chapter file pinned-original-source.tex carries only \input{preamble} on its first line, so no
% preamble text lives inside the chapter itself.  The Stacks Project uses one single shared preamble for
% every chapter; the snapshot stages it identically as sources/stacks-chow/preamble.tex and
% sources/stacks-etale-cohomology/preamble.tex.  The block below is copied mechanically, line for line,
% from sources/stacks-chow/preamble.tex lines 14-18 (the Xy-pic setup) and 154-217 (the native theorem
% declarations together with the mathematical macros).  Nothing was added, removed or redefined; the
% wrapper class and the amsmath/amssymb/amsthm/graphicx base packages are supplied by the assembler.
% Chapter: more-algebra.tex -- More on Algebra.
% For commutative diagrams we use Xy-pic
\usepackage[all]{xy}

% We use 2cell for 2-commutative diagrams.
\xyoption{2cell}

% Theorem environments.
%
\theoremstyle{plain}
\newtheorem{theorem}[subsection]{Theorem}
\newtheorem{proposition}[subsection]{Proposition}
\newtheorem{lemma}[subsection]{Lemma}

\theoremstyle{definition}
\newtheorem{definition}[subsection]{Definition}
\newtheorem{example}[subsection]{Example}
\newtheorem{exercise}[subsection]{Exercise}
\newtheorem{situation}[subsection]{Situation}

\theoremstyle{remark}
\newtheorem{remark}[subsection]{Remark}
\newtheorem{remarks}[subsection]{Remarks}

\numberwithin{equation}{subsection}

% Macros
%
\def\lim{\mathop{\mathrm{lim}}\nolimits}
\def\colim{\mathop{\mathrm{colim}}\nolimits}
\def\Spec{\mathop{\mathrm{Spec}}}
\def\Hom{\mathop{\mathrm{Hom}}\nolimits}
\def\Ext{\mathop{\mathrm{Ext}}\nolimits}
\def\SheafHom{\mathop{\mathcal{H}\!\mathit{om}}\nolimits}
\def\SheafExt{\mathop{\mathcal{E}\!\mathit{xt}}\nolimits}
\def\Sch{\mathit{Sch}}
\def\Mor{\mathop{\mathrm{Mor}}\nolimits}
\def\Ob{\mathop{\mathrm{Ob}}\nolimits}
\def\Sh{\mathop{\mathit{Sh}}\nolimits}
\def\NL{\mathop{N\!L}\nolimits}
\def\CH{\mathop{\mathrm{CH}}\nolimits}
\def\proetale{{pro\text{-}\acute{e}tale}}
\def\etale{{\acute{e}tale}}
\def\QCoh{\mathit{QCoh}}
\def\Ker{\mathop{\mathrm{Ker}}}
\def\Im{\mathop{\mathrm{Im}}}
\def\Coker{\mathop{\mathrm{Coker}}}
\def\Coim{\mathop{\mathrm{Coim}}}

% Boxtimes
%
\DeclareMathSymbol{\boxtimes}{\mathbin}{AMSa}{"02}

%
% Macros for moduli stacks/spaces
%
\def\QCohstack{\mathcal{QC}\!\mathit{oh}}
\def\Cohstack{\mathcal{C}\!\mathit{oh}}
\def\Spacesstack{\mathcal{S}\!\mathit{paces}}
\def\Quotfunctor{\mathrm{Quot}}
\def\Hilbfunctor{\mathrm{Hilb}}
\def\Curvesstack{\mathcal{C}\!\mathit{urves}}
\def\Polarizedstack{\mathcal{P}\!\mathit{olarized}}
\def\Complexesstack{\mathcal{C}\!\mathit{omplexes}}
% \Pic is the operator that assigns to X its picard group, usage \Pic(X)
% \Picardstack_{X/B} denotes the Picard stack of X over B
% \Picardfunctor_{X/B} denotes the Picard functor of X over B
\def\Pic{\mathop{\mathrm{Pic}}\nolimits}
\def\Picardstack{\mathcal{P}\!\mathit{ic}}
\def\Picardfunctor{\mathrm{Pic}}
\def\Deformationcategory{\mathcal{D}\!\mathit{ef}}

% --- appended 2026-10-12 (Chief): the project own definitions of the reference/slogan/history
% display environments, copied mechanically from the shared Stacks preamble lines 8-12
% (\usepackage{verbatim} plus the three \newenvironment lines, which the chapter preamble keeps
% verbatim from the project preamble). Nothing added, removed or redefined: these environments
% carry no mathematical content and the project itself typesets them as comments. Without these
% three lines a counted statement whose own text uses them fails with Environment slogan undefined.
\usepackage{verbatim}
\newenvironment{reference}{\comment}{\endcomment}
\newenvironment{slogan}{\comment}{\endcomment}
\newenvironment{history}{\comment}{\endcomment}
\usepackage{cleveref}
\crefname{originallemma}{original lemma}{original lemmas}
\Crefname{originallemma}{Original lemma}{Original lemmas}
\makeatletter
\newlabel{algebra-lemma-exact-sequence-projective-dimension}{{065S}{source}{}{sourcepointer-1}{}}
\newlabel{algebra-lemma-exact-sequence-projective-dimension@cref}{{[originallemma][0][]065S}{[1][0][]source}{}{}{}}
\makeatother
\crefname{originallemma}{original lemma}{original lemmas}
\Crefname{originallemma}{Original lemma}{Original lemmas}
\makeatletter
\newlabel{algebra-lemma-finite-gl-dim-finite-dim-regular}{{00OE}{source}{}{sourcepointer-2}{}}
\newlabel{algebra-lemma-finite-gl-dim-finite-dim-regular@cref}{{[originallemma][0][]00OE}{[1][0][]source}{}{}{}}
\makeatother
\crefname{originallemma}{original lemma}{original lemmas}
\Crefname{originallemma}{Original lemma}{Original lemmas}
\makeatletter
\newlabel{derived-definition-generators}{{09SJ}{source}{}{sourcepointer-3}{}}
\newlabel{derived-definition-generators@cref}{{[originallemma][0][]09SJ}{[1][0][]source}{}{}{}}
\makeatother
\crefname{originallemma}{original lemma}{original lemmas}
\Crefname{originallemma}{Original lemma}{Original lemmas}
\makeatletter
\newlabel{derived-lemma-ext-2-zero-pre}{{0GM4}{source}{}{sourcepointer-4}{}}
\newlabel{derived-lemma-ext-2-zero-pre@cref}{{[originallemma][0][]0GM4}{[1][0][]source}{}{}{}}
\makeatother
\crefname{originallemma}{original lemma}{original lemmas}
\Crefname{originallemma}{Original lemma}{Original lemmas}
\makeatletter
\newlabel{lemma-regular-perfect}{{19115}{source}{}{sourcepointer-5}{}}
\newlabel{lemma-regular-perfect@cref}{{[originallemma][0][]19115}{[1][0][]source}{}{}{}}
\makeatother
\title{A Noetherian regular ring of finite dimension is a strong generator for its perfect complexes}
\author{The Stacks Project Authors}
\date{stacks-project:more-algebra.tex:local-snapshot-20261012; GNU FDL 1.2 or later}
\begin{document}
\maketitle
\setcounter{subsection}{1}
\noindent\textit{Editorial note.} All delivered mathematical text is the author's own, copied from the pinned Stacks Project chapter file more-algebra.tex; no statement or proof step was rewritten, reordered or removed.\par\medskip
\noindent\textit{Editorial note.} The chapter file pinned-original-source.tex carries only \textbackslash{}input\{preamble\} on its first line, so no preamble text lives inside the chapter. The retained file sources/stacks-more-algebra/retained-preamble.tex is a mechanical, line-for-line copy of the project's single shared preamble as staged in the snapshot; nothing was added, removed or redefined and no package was added.\par\medskip
\noindent\textit{Editorial note.} The source URL is the source-level Stacks Project URL, as required when more than one card is produced from one source; the per-claim tag is carried by claim\_key.\par\medskip
\noindent\textit{Editorial note.} The author writes '(small detail omitted)' once inside the counted proof span, at the lifting of F\textasciicircum{}n -> H\textasciicircum{}n(K) to F\textasciicircum{}n[-n] -> K; that phrase is kept verbatim and the surrounding argument is complete without it. The five references inside the counted proof span are recorded as non-load-bearing navigation pointers.\par\medskip
\noindent\textit{Allowed standard prerequisites (editorial).}\begin{itemize}
\item Perfect objects as bounded complexes with finite cohomology, and strong generation of a triangulated category. invoked at proof lines 20827-20831. Fixes the two notions the statement relates and the generation index that the induction tracks.
\item Finite global dimension of a regular ring and subadditivity of projective dimension in a short exact sequence. invoked at proof lines 20836-20837 and 20875-20877. Supplies the bound d and is what makes the projective dimension of H\textasciicircum{}n(K') drop by one.
\item Vanishing of higher Ext between projective modules and the splitting of a complex into the shifts of its cohomology. used at proof lines 20845-20851. Is the base case of the induction.
\end{itemize}\medskip
\ifdefined\hypertarget\hypertarget{sourcepointer-1}{}\fi
\noindent\textit{Source reference.} Original lemma 065S: focused-seeds/AGP/sources/stacks-snapshot-20261011/tags.txt, source lines 7911--7911. Cross-chapter navigation pointer into the Stacks Project; the standard fact it supplies is applied inside the counted proof span and the argument resting on it stays there.\par\medskip
\ifdefined\hypertarget\hypertarget{sourcepointer-2}{}\fi
\noindent\textit{Source reference.} Original lemma 00OE: focused-seeds/AGP/sources/stacks-snapshot-20261011/tags.txt, source lines 892--892. Cross-chapter navigation pointer into the Stacks Project; the standard fact it supplies is applied inside the counted proof span and the argument resting on it stays there.\par\medskip
\ifdefined\hypertarget\hypertarget{sourcepointer-3}{}\fi
\noindent\textit{Source reference.} Original lemma 09SJ: focused-seeds/AGP/sources/stacks-snapshot-20261011/tags.txt, source lines 12357--12357. Cross-chapter navigation pointer into the Stacks Project; the standard fact it supplies is applied inside the counted proof span and the argument resting on it stays there.\par\medskip
\ifdefined\hypertarget\hypertarget{sourcepointer-4}{}\fi
\noindent\textit{Source reference.} Original lemma 0GM4: focused-seeds/AGP/sources/stacks-snapshot-20261011/tags.txt, source lines 20744--20744. Cross-chapter navigation pointer into the Stacks Project; the standard fact it supplies is applied inside the counted proof span and the argument resting on it stays there.\par\medskip
\ifdefined\hypertarget\hypertarget{sourcepointer-5}{}\fi
\noindent\textit{Source reference.} Original lemma 19115: sources/stacks-more-algebra/pinned-original-source.tex, source lines 19115--19115. Same-chapter navigation pointer; the statement it names is standard for the section and the argument that consumes it stays inside the counted proof span.\par\medskip
\begin{lemma}
\label{lemma-regular-strong-generator}
Let $R$ be a Noetherian regular ring of finite dimension.
Then $R$ is a strong generator for the full subcategory
$D_{perf}(R) \subset D(R)$ of perfect objects.
\end{lemma}

\begin{proof}
We will use that an object $K$ of $D(R)$ is perfect if and only
if $K$ is bounded and has finite cohomology modules, see
Lemma \ref{lemma-regular-perfect}.
Strong generators of triangulated categories are defined in
Derived Categories, Definition \ref{derived-definition-generators}.
Let $d = \dim(R)$.

\medskip\noindent
Let $K \in D_{perf}(R)$. We will show $K \in \langle R \rangle_{d + 1}$.
By Algebra, Lemma \ref{algebra-lemma-finite-gl-dim-finite-dim-regular}
every finite $R$-module has projective dimension $\leq d$.
We will show by induction on $0 \leq i \leq d$ that
if $H^n(K)$ has projective dimension $\leq i$ for all
$n \in \mathbf{Z}$, then $K$ is in $\langle R \rangle_{i + 1}$.

\medskip\noindent
Base case $i = 0$. In this case $H^n(K)$ is a finite $R$-module
of projective dimension $0$. In other words, each cohomology
is a projective $R$-module. Thus $\Ext^i_R(H^n(K), H^m(K)) = 0$
for all $i > 0$ and $m, n \in \mathbf{Z}$. By
Derived Categories, Lemma \ref{derived-lemma-ext-2-zero-pre}
we find that $K$ is isomorphic to the direct sum of the shifts of its
cohomology modules. Since each cohomology module is a finite
projective $R$-module, it is a direct summand of a direct sum
of copies of $R$. Hence by definition we see that $K$
is contained in $\langle R \rangle_1$.

\medskip\noindent
Induction step. Assume the claim holds for $i < d$ and let
$K \in D_{perf}(R)$ have the property that $H^n(K)$ has
projective dimension $\leq i + 1$ for all $n \in \mathbf{Z}$.
Choose $a \leq b$ such that $H^n(K)$ is zero for $n \not \in [a, b]$.
For each $n \in [a, b]$ choose a surjection $F^n \to H^n(K)$
where $F^n$ is a finite free $R$-module. Since $F^n$ is projective,
we can lift $F^n \to H^n(K)$ to a map $F^n[-n] \to K$ in $D(R)$
(small detail omitted). Thus we obtain a morphism
$\bigoplus_{a \leq n \leq b} F^n[-n] \to K$
which is surjective on cohomology modules.
Choose a distinguished triangle
$$
K' \to \bigoplus\nolimits_{a \leq n \leq b} F^n[-n] \to K \to K'[1]
$$
in $D(R)$. Of course, the object $K'$ is bounded and has
finite cohomology modules. The long exact sequence of cohomology
breaks into short exact sequences
$$
0 \to H^n(K') \to F^n \to H^n(K) \to 0
$$
by the choices we made. By
Algebra, Lemma \ref{algebra-lemma-exact-sequence-projective-dimension}
we see that the projective dimension of $H^n(K')$ is $\leq \max(0, i)$.
Thus $K' \in \langle R \rangle_{i + 1}$. By definition this means
that $K$ is in $\langle R \rangle_{i + 1 + 1}$ as desired.
\end{proof}

\end{document}
